\documentclass[11pt]{article}
\usepackage{ap-psych}
\title{Research Methods \& Data: Question Set}
\author{Paul Wang}
\date{November 2026}
\docpart{Part 0 \textperiodcentered\ Question Set}

\begin{document}
\maketitle

\begin{abstract}
Exam-style practice on the skills the redesigned exam tests in every unit: identifying methods and variables, judging what a design can and cannot conclude, reading statistics and graphs, and applying ethical guidelines. Questions are grouped by skill and several are built on shared study descriptions, as on the real exam. A full Article Analysis Question closes the set. All items are original. Allow about 25 minutes for the multiple choice and 25 minutes for the AAQ.
\end{abstract}

\scoringline{MCQ: \blank\ / \qtotal \qquad AAQ: \blank\ / 7 \qquad Target: 18+ MCQ and 5+ AAQ \qquad Pace: 72 s per MCQ}

\section{Identifying methods and variables}

\begin{stimulus}
\textbf{Study A.} Researchers recruited 120 adults and randomly assigned them to walk outdoors in a park or along a busy street for 50 minutes. Before and after the walk, participants completed the Rumination Scale (scores 0--60). The park group's scores dropped by an average of 6 points; the street group's scores did not change.
\end{stimulus}

\begin{q}{2}{RM-2 Variables}{Walking setting is the manipulated variable.}
In Study A, the independent variable is
\ch{the Rumination Scale score}{the walking environment (park vs. street)}{the 50-minute walk length}{the participants' ages}
\end{q}

\begin{q}{4}{RM-2 Variables}{The operational definition of the DV is the specific measure used.}
In Study A, rumination is operationally defined as
\ch{repetitive negative thinking}{the time spent outdoors}{the difference between the two groups}{the score on the Rumination Scale before and after the walk}
\end{q}

\begin{q}{1}{RM-2 Methods}{Random assignment + manipulated IV permits a cause-and-effect conclusion.}
Which conclusion is best supported by Study A?
\ch{Walking in a park caused a reduction in rumination compared with walking on a busy street.}{People who ruminate less prefer parks.}{Rumination is correlated with age.}{Nature walks cure depression.}
\end{q}

\begin{q}{3}{RM-2 Methods}{Pre-existing groups that were not assigned \(\rightarrow\) correlational/quasi; no causal claim.}
A psychologist compares the anxiety scores of students who already play a varsity sport with students who do not. This design is best described as
\ch{an experiment, because there are two groups}{a case study, because it focuses on athletes}{correlational, because the groups were not randomly assigned}{longitudinal, because anxiety is measured}
\end{q}

\begin{q}{2}{RM-2 Methods}{Combining many prior studies statistically = meta-analysis.}
A researcher statistically combines the results of 48 published studies on mindfulness and blood pressure to estimate an overall effect. This is
\ch{a longitudinal study}{a meta-analysis}{naturalistic observation}{a cross-sectional study}
\end{q}

\begin{q}{4}{RM-2 Methods}{Following the same children over years is longitudinal.}
To study how self-control changes, researchers test the same 300 children at ages 4, 8, 12, and 16. A drawback of this design is
\ch{cohort effects}{the lack of a dependent variable}{that it cannot measure change within individuals}{attrition as participants drop out over time}
\end{q}

\section{Validity, bias, and controls}

\begin{q}{1}{RM-3 Confound}{Different instructors differ in more than the IV, so instructor is a confound.}
A teacher tests a new study app by giving it to her morning class and not to her afternoon class, which another teacher instructs. The most serious problem is that
\ch{the classes differ in teacher and time of day, which are confounding variables}{the sample is too large}{the study uses random assignment}{there is no operational definition of ``app''}
\end{q}

\begin{q}{3}{RM-3 Bias}{Respondents shade answers to look good: social desirability bias.}
On a survey, many respondents underreport how often they skip class. This most likely reflects
\ch{the placebo effect}{experimenter bias}{social desirability bias}{random assignment failure}
\end{q}

\begin{q}{2}{RM-3 Controls}{A placebo group separates the drug's effect from expectation.}
Why do drug trials give a control group an inert sugar pill?
\ch{To increase the sample size}{To separate the drug's effect from the effect of expecting treatment}{To ensure participants are randomly sampled}{To make the study longitudinal}
\end{q}

\begin{q}{4}{RM-3 Bias}{Observers who know the hypothesis may see what they expect.}
Observers who know which children watched a violent video rate those children as more aggressive on the playground. The best fix is to
\ch{use a larger sample}{replace the video with a survey}{use only one observer}{keep observers unaware of each child's condition}
\end{q}

\begin{q}{1}{RM-3 Generalizability}{Random sampling from the target population supports generalization.}
Which change would most improve the generalizability of a study run with 40 volunteers from one gym?
\ch{Randomly sampling adults from the city's population}{Making the study double-blind}{Adding a placebo condition}{Running the study for a longer time}
\end{q}

\section{Reading statistics and data}

\begin{stimulus}
\textbf{Study B.} A survey of 500 high-school students measured daily hours of phone use after 10 p.m.\ and average nightly sleep. The correlation was \(r = -.58\). Students in the top quartile of phone use averaged 6.1 hours of sleep (SD = 0.9); students in the bottom quartile averaged 7.4 hours (SD = 0.8).
\end{stimulus}

\begin{q}{3}{RM-4 Correlation}{Negative r: as one variable rises, the other falls.}
The correlation in Study B indicates that
\ch{late phone use causes students to sleep less}{students who sleep less use their phones less}{students who use phones more at night tend to sleep less}{there is no meaningful relationship}
\end{q}

\begin{q}{2}{RM-4 Causation}{Correlation leaves third variables and reverse causation open.}
Which is the best reason Study B cannot show that phone use reduces sleep?
\ch{The sample is too small.}{A third variable, such as anxiety, could increase both phone use and sleeplessness.}{The correlation is positive.}{Sleep was not operationally defined.}
\end{q}

\begin{q}{4}{RM-4 SD}{Similar SDs mean similar spread; groups differ in mean, not variability.}
What do the standard deviations in Study B show?
\ch{The top-quartile group slept far less consistently.}{The bottom-quartile group had a wider range of phone use.}{The difference in means is statistically significant.}{Sleep scores were about equally spread out in both groups.}
\end{q}

\begin{q}{1}{RM-4 Scatterplot}{r = −.58: points trend downward with moderate scatter.}
A scatterplot of Study B's data would most likely show
\ch{points sloping downward from left to right with moderate scatter}{points forming a tight upward line}{points scattered with no pattern}{a perfect horizontal line}
\end{q}

\begin{q}{3}{RM-4 Central tendency}{Mode = most frequent score (4 appears three times).}
Scores on a 10-point stress scale were 2, 4, 4, 4, 5, 7, 9. What is the mode?
\chh{5}{4.5}{4}{7}
\end{q}

\begin{q}{2}{RM-4 Skew}{An easy test bunches scores high with a tail to the left: negative skew, mean < median.}
Most students scored 90--100 on an easy quiz, with a few scores near 40. The distribution is
\ch{positively skewed, with the mean above the median}{negatively skewed, with the mean below the median}{normal, with mean equal to median}{bimodal}
\end{q}

\begin{q}{1}{RM-4 Effect size}{Significance with a tiny difference in a huge sample = small effect size.}
With 20{,}000 participants, a memory training program improves recall by 0.3 words (\(p < .001\)). The best interpretation is that the result is
\ch{statistically significant but small in practical size}{not statistically significant}{proof that the program works for everyone}{caused by sampling bias}
\end{q}

\section{Ethics}

\begin{q}{4}{RM-5 Ethics}{Minors need parental consent plus their own assent.}
Before 12-year-olds join a study, the researchers obtain permission from parents and agreement from each child. The child's agreement is called
\ch{debriefing}{confidentiality}{deception}{assent}
\end{q}

\begin{q}{3}{RM-5 Ethics}{The IRB reviews proposals for risk to participants before research begins.}
Which body reviews a research proposal to ensure participants are protected before the study begins?
\ch{The American Psychiatric Association}{A meta-analysis committee}{An Institutional Review Board}{The participants' teachers}
\end{q}

\begin{q}{2}{RM-5 Ethics}{Participants may leave at any point without penalty.}
Midway through a study, a participant feels uncomfortable and leaves; she still receives course credit. Which guideline is honored?
\ch{Debriefing}{The right to withdraw}{Deception}{Random assignment}
\end{q}

\section{Article Analysis Question}

\begin{frq}{Question 1 \textperiodcentered\ AAQ \textperiodcentered\ 7 points \textperiodcentered\ 25 minutes}
\small
\textbf{Source (summarized).} Okafor, L., \& Brandt, S. (2024). Retrieval practice and exam performance in secondary students. \emph{Journal of Applied Learning}, 18, 211--229.

\textbf{Participants.} 186 students (ages 15--17) from three public high schools in the Midwestern United States. Parents gave written consent and students gave assent. Students could stop at any time; data were stored under ID numbers.

\textbf{Method.} Students were randomly assigned to one of two study conditions for a unit on cell biology. The \emph{retrieval} group answered 10 low-stakes quiz questions at the end of each lesson. The \emph{restudy} group reread their notes for the same amount of time. One week after the unit, all students took the same 40-question exam, graded by teachers who did not know each student's condition.

\textbf{Results.} The retrieval group averaged 31.2 correct (SD = 4.1); the restudy group averaged 26.5 (SD = 4.4). The difference was statistically significant (\(p < .01\)).

\medskip
\begin{enumerate}[label=(\Alph*),leftmargin=2em,itemsep=3pt]
\item Identify the research method used in the study.
\item State the operational definition of exam performance.
\item Describe what the reported \(p\) value indicates about the results.
\item Identify at least one ethical guideline applied by the researchers.
\item Explain the extent to which the findings may or may not be generalizable, using specific information from the study.
\item Explain how at least one finding supports or refutes the claim that ``testing yourself is a more effective way to encode information into long-term memory than rereading.''
\end{enumerate}
\end{frq}

\makekey

\subsection*{AAQ scoring guide (7 points)}
\small
\begin{compare}{Part}{Earns the point}{Does not earn the point}
A (1) & ``Experiment.'' & ``Correlational study''; ``case study.'' \\
B (1) & Number correct out of 40 on the same exam given one week after the unit. & ``How well students did''; ``their grade'' with no measure. \\
C (1) & The difference between the retrieval and restudy groups is unlikely to be due to chance (less than a 1\% probability). & ``The effect was large''; ``1\% of students improved.'' \\
D (1) & Informed consent (parental consent and student assent), right to withdraw, \emph{or} confidentiality via ID numbers, tied to what the study did. & Naming ``debriefing'' or ``protection from harm'' with no link to the source. \\
E (1) & Limited: only 186 students aged 15--17 from three Midwestern U.S. schools, studying one biology unit, so results may not extend to adults, other regions, or other subjects. \emph{Or} supported: random assignment and blind grading make the effect credible across similar classrooms. & ``It's only one study''; ``the sample is small'' with no specifics. \\
F (2) & 1 pt: states the finding supports the claim and cites a result (31.2 vs.\ 26.5 correct, \(p < .01\)). 1 pt: explains with a concept, e.g., the \emph{testing effect}: effortful retrieval strengthens memory traces and improves later recall more than passive rereading. & Restating the claim; describing the testing effect without tying it to the study's numbers. \\
\end{compare}

\end{document}
